% Generated draft: requires mathematical and source review.
\documentclass[11pt]{article}
\usepackage[margin=25mm]{geometry}
\usepackage{fontspec}
\setmainfont{Latin Modern Roman}
\usepackage{amsmath,amssymb,mathtools,unicode-math}
\setmathfont{Latin Modern Math}
\usepackage{enumitem,needspace,xurl,hyperref}
\setlength{\parindent}{0pt}
\setlength{\parskip}{5pt}
\setlength{\emergencystretch}{3em}
\newcommand{\N}{\mathbb{N}}
\newcommand{\Z}{\mathbb{Z}}
\newcommand{\Q}{\mathbb{Q}}
\newcommand{\R}{\mathbb{R}}
\newcommand{\C}{\mathbb{C}}
\newcommand{\F}{\mathbb{F}}
\newcommand{\A}{\mathbb{A}}
\newcommand{\PP}{\mathbb{P}}
\newcommand{\E}{\mathbb{E}}
\newcommand{\eps}{\varepsilon}
\newcommand{\dd}{\,\mathrm{d}}
\begin{document}
\section*{UnsolvedMath ID: 20002387}
\textbf{Source problem code:} catalogue entry
\subsection{Definitions and mathematical statement}
For a homogeneous polynomial \(p:\R^d\to\R\), hyperbolicity with respect to \(e\) means \(p(e)>0\) and that \(t\mapsto p(x-te)\) has only real roots for every \(x\). Its closed hyperbolicity cone \(\Lambda_+(p,e)\) is the closure of the connected component of \(\{p>0\}\) containing \(e\). A linear matrix inequality (LMI) has the form \(A_0+\sum_j y_jA_j\succeq0\), with specified real symmetric matrices.

The literal source defines a hyperbolic feasibility problem by polynomials \(p_1,\ldots,p_m\) and raw inequalities \(p_i(x)\ge0\), without choosing hyperbolicity components. Given an ellipsoid \(E_0\subseteq\R^d\) of volume \(R>1\) and \(\varepsilon\in(0,1)\), its proposed task is to output either
\[
 x\in E_0\cap\bigcap_{i=1}^m\{p_i\ge0\},
\]
or the correct assertion that this intersection contains no Euclidean ball of volume at least \(\varepsilon\). The second output is a weak infeasibility certificate and does not say the feasible set is empty. The proposed size is
\[
 s=\max\{m,\log(1/\varepsilon),\log R\}.
\]
Open Problem 1.10.3 asks whether every such problem can be explicitly converted to an LMI, solvable to the stated accuracy in \(\operatorname{poly}(s)\) Blum--Shub--Smale real operations, whose solution supplies the required original output. This is the source's stated reduction requirement, despite its use of the word ``equivalent.'' The OCR control character is interpreted as \(\varepsilon\); a displayed \(\operatorname{size}(\mathrm{HLP})\) appears to be a typographical variant of the HFP size.

The aid identifies two separate defects. First, \(s\) omits \(d\), the polynomial representation, and the ellipsoid data. Under a convention charging explicit coordinate output, a fixed-\((m,R,\varepsilon)\) family with unbounded \(d\) requires \(\Omega(d)\) output while \(s\) is constant. One can take the linear hyperbolic constraint \(p(x)=x_1\) and a ball \(E_0\) of fixed volume \(R\), translated wholly into \(\{x_1>0\}\). Its feasible set contains that ball, so the no-ball output is false for fixed \(\varepsilon<R\). This obstruction depends on the stated explicit-output cost convention and refutes that size bound, not every repaired reduction question.

Second, raw hyperbolic inequalities need not define a convex set. For \(p(x,y)=xy\), hyperbolic with respect to \((1,1)\), the set \(\{xy\ge0\}\) is nonconvex. Every projection of a spectrahedron is convex. Hence this literal set cannot be exactly such a projection. This does not rule out an arbitrary computational reduction that does not preserve the entire feasible set.

The supplied conditional repair includes complete input encoding, chosen cone directions, appropriate conditioning and accuracy data, and explicit spectrahedral lifts
\[
 \Lambda_+(p_i,e_i)=\{x:\exists z_i\ L_i(x,z_i)\succeq0\}
\]
whose total encoded size is polynomial in the repaired input size. If the encoded ellipsoid is \(E_0=\{x:\|B(x-c)\|_2\le1\}\), its LMI is
\[
 L_E(x)=\begin{pmatrix}1&(B(x-c))^{\mathsf T}\\ B(x-c)&I_d\end{pmatrix}\succeq0.
\]
Then
\[
 \operatorname{diag}(L_1(x,z_1),\ldots,L_m(x,z_m),L_E(x))\succeq0
\]
projects exactly onto \(E_0\cap\bigcap_i\Lambda_+(p_i,e_i)\), with polynomial total size under the supplied-lift assumption. This is an exact representation conditional on having those lifts. It neither establishes their existence for every hyperbolicity cone nor, by itself, proves the source's volume-sensitive accuracy and running-time guarantee: conditioning and recovery of a suitable approximate solution require separate quantitative control. The final extracted phrase ``A problem analogous to'' is truncated and is not completed here.
\subsection{Short English statement}
The printed feasibility problem omits essential input size and cone choices. The aid separates an explicit-output size obstruction from a nonconvexity obstruction to exact LMI projections. A block LMI represents the repaired cone intersection only when polynomial-size lifts and the needed numerical data are supplied.
\subsection{Sources}
\begin{itemize}
\item Local registry, permanent UnsolvedMath ID \(20002387\), exact public code \texttt{catalogue entry}; canonical location \texttt{aim-optimization-notes.json}, \texttt{notes[2]}. The AIM domain is Optimization; the registry category is Analysis.
\item AIM, \emph{Theory and Algorithms of Linear Matrix Inequalities}, Open Problem 1.10.3 and Definition 1.10.2, printed pages 10--11, PDF pages 11--12, version 12 March 2006; supplied reference \url{https://aimath.org/WWN/matrixineq/matrixineq.pdf}. No proposer, original remarks, or literature entry is supplied. The statement is marked \texttt{reconstructed\_unverified}.
\item Supplied research attempt 1: \texttt{invalid\_statement}, run status \texttt{unknown}, rigor label \texttt{proved}, low novelty confidence; registry status \texttt{open}, classification \texttt{OPEN-TRIAGE}, created and updated 14 August 2026. The two obstructions and conditional lift construction have no additional primary citations in the record.
\item The research aid is machine generated, not peer reviewed, and explicitly requires independent verification before citation.
\end{itemize}
\end{document}
